% page 505 of the facsimile (image p-504 of the scan)
% block 175.3 x 242.0 mm ; left margin 26.7 mm
\pgc{5.080mm}{0.000mm}{
\l{\cel{50}497}
\l{}
\l{\sou{revokar}. (\sou{trans.}) (\sou{pro}). Retrotirar (lu) de lua ofico (pro neka-}
\l{\cel{3}pableso, pro kulpo profesionala, e c.) - EFIS.}
\l{}
\l{\sou{revolta}r (\sou{netrans., kontre, pri, pro)}. Refuzar la obedio di la im-}
\l{\cel{3}peri de la autoritatozi legitima, regnanta. - (\sou{anke metaf}.)}
\l{\cel{3}DEFIRS.}
\l{}
\l{\sou{revolucionar}. (\sou{netrans.})\cel{2}(\sou{aludante la populo di stato)}.- Agar}
\l{\cel{3}violente la chanjo di (sua) guvernistaro. - DEFIRS.}
\l{}
\l{\sou{revolvero}. Pafarmo portebla, pistoleto qua posibligas pafar plura-}
\l{\cel{3}foye sen interrupto, per provizesir ye plura kulati quin onu}
\l{\cel{3}povas charjar anticipe, e quin resorto igas pasar singlope avan}
\l{\cel{3}la tubo o la tubi di la armo. - DEFIRS.}
\l{}
\l{\sou{revuar}.\cel{2}(\sou{trans.}) I. (\sou{armeo}) Manovro, defilo, di la armo-trupi,}
\l{\cel{3}di qua la skopo esas certigar su (igar su +\sou{certena})pri la ko-}
\l{\cel{3}rekteso di lia su-teno,.instrukteso, precizeso di movi e gesti.}
\l{\cel{3}- II. Explorar la aktualaji en edituro periodala. - DEFIS.}
\l{}
\l{\sou{revulsar}. (\sou{trans}.) (\sou{patol.}) Igar inflamantajo, obstruktantajo, e c.}
\l{\cel{3}deviacar de la loko malada, per atraktar lu ad altra parto di la}
\l{\cel{3}korpo ube lu ne esas danjeroza. - EF.}
\l{}
\l{\sou{reza}. I. (\sou{pri pili, hari, e}\cel{1}c.) Tondita tale ke li ne salias de la}
\l{\cel{3}surfaco cirkondanta. - II. Di qua la surfaco esas sen saliaji}
\l{\cel{3}nek sulki. - III. (pri recipiento e lua kontenajo solida). Qua}
\l{\cel{3}ne salias super la bordo o bordi. - EF.}
\l{}
\l{\sou{rezedo}. (\sou{bot.)}\cel{2}Planto herbacea, kun flori odoranta, folii dispo-}
\l{\cel{3}zita sur la stipo intervaloze, quan onu kultivas en la gardeni.}
\l{\cel{3}- L.\cel{1}\sou{reseda}. DEFIRS.}
\l{}
\l{\sou{rezervar}. (\sou{trans.}) I. Retenar, gardar, ula-destine. - II. Ne lasar}
\l{\cel{3}sua impresesi, sua sentimenti, konocesar da kunhomi. - (\sou{Ca-kaze,}}
\l{\cel{3}\sou{uzesas nur ye la formo \textquotedbl{}esar rezervita\textquotedbl{}}.) DEFIRS.}
\l{}
\l{\sou{rezidar}. (\sou{netrans., en , che}). I. Esar establisita nun en ta o ca}
\l{\cel{3}loko. Permanar en la loko ube onu havas sua ofico. - II. Havar}
\l{\cel{3}sua +siejo (en, che). - DEFIS.}
\l{}
\l{\sou{rezidento}.\cel{2}Agento diplomacala qua rezidas apud la guvernistaro}
\l{\cel{3}di mikra lando, kun grado infra ye olta di ambasadisto o di}
\l{\cel{3}plenipotenciero. - DEFIS.}
\l{}
\l{\sou{reziduo}. La materio qua restas pos operaco kemiala, pos manipulo}
\l{\cel{3}industriala ; ta parto di la dishi qua ne manjesis dum la repasto;}
\l{\cel{3}la restajo di objekto ruptita acidente. - DEFIS.}
\l{}
\l{\sou{rezignar}. (\sou{netrans., pri}) Aceptar volunte desfeliceso, domajo, perdo,}
\l{\cel{3}quan onu subisis e, generale, ne povis evitar.- DEFIS.}
\l{}
\l{\sou{rezino}. Materio inflamebla qua defluas, naturale o per incizo,}
\l{\cel{3}de la pino, pinastro, abieto. - EFIS.}
}
